\subsection{Bilinear Form Types and Congruence Invariants}\label{subsec:metric-vector-spaces-and-bilinear-forms}

% **Role:** arbitrary-field metric geometry.
%
% Symmetric, alternating, and skew-symmetric bilinear forms; rank, discriminant,
% and canonical forms under congruence.

Metric vector spaces retain the algebraic structure of inner-product geometry without imposing positivity. Throughout this subsection, $\mathbb{F}$ is an arbitrary field and all vector spaces are finite-dimensional unless stated otherwise.\\
General bilinear pairings, their left and right radicals, their representing matrices, and the corresponding basis-change law were established earlier. Here the two input spaces coincide, and a fixed bilinear form $B$ is written $B(\mathbf{x},\mathbf{y})=\innerprod{\mathbf{x},\mathbf{y}}$.

\begin{definition} \label{def:metric-vector-space-basic-form-types} \hlt{Symmetric, Skew-Symmetric, and Alternating Forms}\\
Let $\mathcal{V}$ be a vector space over $\mathbb{F}$ and let $B:\mathcal{V}\times\mathcal{V}\to\mathbb{F}$ be bilinear.
\begin{enumerate}[label=(\roman*)]
\setlength{\itemsep}{0pt}
\item $B$ is a \hlt{bilinear form} on $\mathcal{V}$; when fixed, write $B(\mathbf{x},\mathbf{y})=\innerprod{\mathbf{x},\mathbf{y}}$.
\item $\operatorname{Sym}(B)$ means $(\forall\,\mathbf{x},\mathbf{y}\in\mathcal{V})\,B(\mathbf{x},\mathbf{y})=B(\mathbf{y},\mathbf{x})$; then $B$ is \hlt{symmetric}.
\item $\operatorname{Skew}(B)$ means $(\forall\,\mathbf{x},\mathbf{y}\in\mathcal{V})\,B(\mathbf{x},\mathbf{y})=-B(\mathbf{y},\mathbf{x})$; then $B$ is \hlt{skew-symmetric}.
\item $\operatorname{Alt}(B)$ means $(\forall\,\mathbf{x}\in\mathcal{V})\,B(\mathbf{x},\mathbf{x})=0$; then $B$ is \hlt{alternating}.
\item A \hlt{metric vector space} is a pair $(\mathcal{V},B)$ such that $\operatorname{Sym}(B)$, $\operatorname{Skew}(B)$, or $\operatorname{Alt}(B)$ holds.\\
The symmetric and alternating cases are called \hlt{orthogonal geometry} and \hlt{symplectic geometry}, respectively.
\end{enumerate}
\end{definition}

\begin{proposition} \label{prop:form-type-relations-by-characteristic} \hlt{Characteristic Controls the Form Types}\\
Let $B(\mathbf{x},\mathbf{y})=\innerprod{\mathbf{x},\mathbf{y}}$ be a bilinear form on $\mathcal{V}$.
\begin{enumerate}[label=(\roman*)]
\setlength{\itemsep}{0pt}
\item $\operatorname{Alt}(B)\Rightarrow\operatorname{Skew}(B)$.
\item If $\operatorname{char}(\mathbb{F})\neq2$, then $\operatorname{Sym}(B)\wedge\operatorname{Alt}(B)\Rightarrow B=0$.
\item If $\operatorname{char}(\mathbb{F})\neq2$, then $\operatorname{Alt}(B)\Longleftrightarrow\operatorname{Skew}(B)$.
\item If $B$ is symmetric, skew-symmetric, or alternating, then $\operatorname{rad}_{L}(B)=\operatorname{rad}_{R}(B)$.
\item If $\operatorname{char}(\mathbb{F})=2$, then $\operatorname{Sym}(B)\Longleftrightarrow\operatorname{Skew}(B)$.
\item If $\operatorname{char}(\mathbb{F})=2$, then $\operatorname{Alt}(B)\Rightarrow\operatorname{Sym}(B)$.
\end{enumerate}
\begin{proof}
\begin{enumerate}[label=(\roman*)]
\setlength{\itemsep}{0pt}
\item[$(i)$] Assume $\operatorname{Alt}(B)$. By bilinearity and alternation, $(\forall\,\mathbf{x},\mathbf{y}\in\mathcal{V})\,B(\mathbf{x},\mathbf{y})+B(\mathbf{y},\mathbf{x})=B(\mathbf{x}+\mathbf{y},\mathbf{x}+\mathbf{y})=0$. Hence $\operatorname{Skew}(B)$.
\item[$(ii)$] Assume $\operatorname{char}(\mathbb{F})\neq2$, $\operatorname{Sym}(B)$, and $\operatorname{Alt}(B)$. By (i), $B$ is also skew-symmetric.\\
Thus $(\forall\,\mathbf{x},\mathbf{y}\in\mathcal{V})\,2B(\mathbf{x},\mathbf{y})=0$. Since $2\neq0$ in $\mathbb{F}$, $(\forall\,\mathbf{x},\mathbf{y}\in\mathcal{V})\,B(\mathbf{x},\mathbf{y})=0$, and therefore $B=0$.
\item[$(iii)$] ($\Rightarrow$) If $\operatorname{Alt}(B)$, then $\operatorname{Skew}(B)$ by (i).
\item[$(iii)$] ($\Leftarrow$) Assume $\operatorname{char}(\mathbb{F})\neq2$ and $\operatorname{Skew}(B)$. Then $(\forall\,\mathbf{x}\in\mathcal{V})\,2B(\mathbf{x},\mathbf{x})=0$. Since $2\neq0$ in $\mathbb{F}$, $(\forall\,\mathbf{x}\in\mathcal{V})\,B(\mathbf{x},\mathbf{x})=0$, so $\operatorname{Alt}(B)$.
\item[$(iv)$] If $B$ is symmetric or skew-symmetric, then $(\forall\,\mathbf{x},\mathbf{y}\in\mathcal{V})\,[B(\mathbf{x},\mathbf{y})=0\Longleftrightarrow B(\mathbf{y},\mathbf{x})=0]$, so $\operatorname{rad}_{L}(B)=\operatorname{rad}_{R}(B)$.\\
If $B$ is alternating, the same conclusion holds because $B$ is skew-symmetric by (i).
\item[$(v)$] ($\Rightarrow$) Assume $\operatorname{char}(\mathbb{F})=2$ and $\operatorname{Sym}(B)$. Then $(\forall\,\mathbf{x},\mathbf{y}\in\mathcal{V})\,B(\mathbf{x},\mathbf{y})=B(\mathbf{y},\mathbf{x})=-B(\mathbf{y},\mathbf{x})$, so $\operatorname{Skew}(B)$.
\item[$(v)$] ($\Leftarrow$) Assume $\operatorname{char}(\mathbb{F})=2$ and $\operatorname{Skew}(B)$. Then $(\forall\,\mathbf{x},\mathbf{y}\in\mathcal{V})\,B(\mathbf{x},\mathbf{y})=-B(\mathbf{y},\mathbf{x})=B(\mathbf{y},\mathbf{x})$, so $\operatorname{Sym}(B)$.
\item[$(vi)$] Assume $\operatorname{char}(\mathbb{F})=2$ and $\operatorname{Alt}(B)$. Then $\operatorname{Skew}(B)$ by (i), and hence $\operatorname{Sym}(B)$ by (v).
\end{enumerate}
\hlt{Structural Commentary}: Items (ii), (iii), (v), and (vi) isolate the characteristic-dependent relations among the form types.\\
\hlt{Key Insight}: Expanding $B(\mathbf{x}+\mathbf{y},\mathbf{x}+\mathbf{y})=0$ converts alternation into the skew-symmetry identity without dividing by $2$.
\end{proof}
\end{proposition}

Henceforth, $B:\mathcal{V}\times\mathcal{V}\to\mathbb{F}$ is usually symmetric, skew-symmetric, or alternating. By Proposition \ref{prop:form-type-relations-by-characteristic}(iv), its left and right radicals agree, so a single radical and a single orthogonality notation suffice.

\begin{definition} \label{def:matrix-rank-and-discriminant-of-form} \hlt{Matrix, Rank, and Discriminant of a Form}\\
Let $\beta=(\mathbf{v}_1,\ldots,\mathbf{v}_n)$ be an ordered basis of a metric vector space $\mathcal{V}$.
\begin{enumerate}[label=(\roman*)]
\setlength{\itemsep}{0pt}
\item The \hlt{matrix of the form} in $\beta$ is $M_{\beta}\defeq[\innerprod{\mathbf{v}_i,\mathbf{v}_j}]_{i,j}$.
\item The \hlt{rank} of the form is $\operatorname{rank}(M_{\beta})$.
\item The \hlt{discriminant set} is $\Delta(\mathcal{V})\defeq\{r^2\det(M_{\beta}):r\in\mathbb{F}^{\times}\}$.
\item If $\det(M_{\beta})\neq0$, then $\Delta(\mathcal{V})=\det(M_\beta)(\mathbb{F}^{\times})^2$ is the square class of $\det(M_\beta)$ in $\mathbb{F}^{\times}/(\mathbb{F}^{\times})^2$.
\end{enumerate}
\end{definition}

\begin{proposition} \label{prop:metric-form-matrix-recovery-and-congruence} \hlt{Matrices and Invariants under Congruence}\\
Let $n=\dim(\mathcal{V})$, let $\beta,\gamma$ be ordered bases of $\mathcal{V}$, and put $C=C_{\beta\leftarrow\gamma}$.
\begin{enumerate}[label=(\roman*)]
\setlength{\itemsep}{0pt}
\item $\innerprod{\mathbf{x},\mathbf{y}}=[\mathbf{x}]_{\beta}^{t}M_{\beta}[\mathbf{y}]_{\beta}$.
\item $M_{\gamma}=C^tM_{\beta}C$.
\item If $A=M_{\beta}$ and $N\in\operatorname{Mat}(n,n,\mathbb{F})$, then $(\exists\,\delta\text{ an ordered basis of }\mathcal{V})\,N=M_{\delta}\Longleftrightarrow(\exists\,P\in\operatorname{Mat}(n,n,\mathbb{F}),\ P\text{ invertible})\,N=P^tAP$.
\item $\operatorname{rank}(M_{\gamma})=\operatorname{rank}(M_{\beta})$.
\item $\Delta(\mathcal{V})$ is independent of the ordered basis used in its definition.
\end{enumerate}
\begin{proof}
\begin{enumerate}[label=(\roman*)]
\setlength{\itemsep}{0pt}
\item[$(i)$] By Proposition \ref{prop:bilinear-coordinate-evaluation}, $\innerprod{\mathbf{x},\mathbf{y}}=[\mathbf{x}]_{\beta}^{t}M_{\beta}[\mathbf{y}]_{\beta}$.
\item[$(ii)$] By Proposition \ref{prop:square-bilinear-form-change-of-basis-congruence}(i), $M_{\gamma}=C^tM_{\beta}C$.
\item[$(iii)$] ($\Rightarrow$) If $N=M_{\delta}$ for an ordered basis $\delta$, then by (ii), $N=C_{\beta\leftarrow\delta}^tAC_{\beta\leftarrow\delta}$, where $C_{\beta\leftarrow\delta}$ is invertible by Theorem \ref{thm:change-of-basis-matrices}(i).
\item[$(iii)$] ($\Leftarrow$) Suppose $N=P^tAP$ with $P=[p_{ij}]$ invertible, and write $\beta=(\mathbf{v}_1,\ldots,\mathbf{v}_n)$. Define $\delta=(\mathbf{w}_1,\ldots,\mathbf{w}_n)$ by $\mathbf{w}_j=\sum\limits_i p_{ij}\mathbf{v}_i$.\\
The columns of $P$ form a basis of $\mathbb{F}^n$ by Theorem \ref{thm:invertible-matrix-columns-basis}. The inverse coordinate map $[\cdot]_{\beta}^{-1}:\mathbb{F}^n\to\mathcal{V}$ is an isomorphism by Proposition \ref{prop:coordinate-map-isomorphism}(iii), so $\delta$ is a basis by Proposition \ref{prop:isomorphisms-preserve-bases}. By construction, $C_{\beta\leftarrow\delta}=P$. Hence $M_{\delta}=P^tAP=N$ by (ii).
\item[$(iv)$] By Theorem \ref{thm:change-of-basis-matrices}(i), $C$ is invertible. By Corollary \ref{cor:transpose-invertible-iff}, $C^t$ is invertible. Hence $\operatorname{rank}(M_{\gamma})=\operatorname{rank}(M_{\beta})$ by (ii), Theorem \ref{thm:full-rank-products-preserve-rank}(ii).
\item[$(v)$] By (ii) and the determinant identities, $\det(M_{\gamma})=\det(C)^2\det(M_{\beta})$. Since $C$ is invertible, put $c\defeq\det(C)\in\mathbb{F}^{\times}$.\\
The bijection $r\mapsto rc$ of $\mathbb{F}^{\times}$ identifies $\{r^2\det(M_{\gamma}):r\in\mathbb{F}^{\times}\}$ with $\{s^2\det(M_{\beta}):s\in\mathbb{F}^{\times}\}$, proving basis independence.
\end{enumerate}
\hlt{Structural Commentary}: Items (ii)--(v) identify congruence with basis change and establish invariance of rank and discriminant.\\
\hlt{Key Insight}: Because basis change acts in both arguments, $M\mapsto P^tMP$; the determinant consequently changes by a nonzero square.
\end{proof}
\end{proposition}
